% 苏布拉马尼扬·钱德拉塞卡（综述）
% license CCBYSA3
% type Wiki

本文根据 CC-BY-SA 协议转载翻译自维基百科\href{https://en.wikipedia.org/wiki/Subrahmanyan_Chandrasekhar}{相关文章}

\begin{figure}[ht]
\centering
\includegraphics[width=6cm]{./figures/4962eacafc5afbdc.png}
\caption{} \label{fig_Sbl_1}
\end{figure}
苏布拉马尼扬·钱德拉塞卡（Subrahmanyan Chandrasekhar，/ˌtʃəndrəˈʃeɪkər/，音译：昌德拉谢卡；[4] 泰米尔语：சுப்பிரமணியன் சந்திரசேகர்，罗马化：Cuppiramaṇiyaṉ Cantiracēkar；1910年10月19日－1995年8月21日）[5] 是一位印裔美国理论物理学家，在恒星结构、恒星演化以及黑洞等领域作出了卓越贡献。他在职业生涯中也曾将相当一部分时间投入于流体力学研究，特别是稳定性与湍流理论，并在这些领域取得了重要成果。钱德拉塞卡与威廉·A·福勒因在“与恒星结构与演化有关的重要物理过程的理论研究”方面的贡献，共同获得1983年诺贝尔物理学奖。他对恒星演化的数学处理为现代天体物理学提供了众多关于大质量恒星及黑洞晚期演化阶段的理论模型。[6][7]许多重要的科学概念、机构与发明均以他命名，包括著名的钱德拉塞卡极限与钱德拉X射线天文台。[8]

钱德拉塞卡出生于英属印度晚期（英属拉杰时期），一生研究范围极广，涉及恒星结构、白矮星、恒星动力学、随机过程、辐射输运、氢负离子的量子理论、流体力学与磁流体稳定性、湍流、椭球平衡体的平衡与稳定性、广义相对论、黑洞的数学理论以及引力波碰撞理论等多个物理领域。[9]在剑桥大学期间，他提出了一个解释白矮星结构的理论模型，首次将电子简并物质中电子速度导致的相对论性质量变化纳入计算。通过此理论，钱德拉塞卡得出结论：白矮星的质量不能超过太阳质量的1.44倍——这一极限即著名的钱德拉塞卡极限。钱德拉塞卡还修正了扬·奥尔特等人最早提出的恒星动力学模型，将银河系中恒星围绕银心旋转时受到的引力场涨落影响纳入考量。他为这一复杂的动力学问题给出了数学解——涉及二十个偏微分方程，并由此提出了一个新的物理量：动力摩擦。这种效应一方面会减缓恒星的运动速度，另一方面又有助于维持星团的整体稳定。钱德拉塞卡进一步将分析扩展至星际介质，证明银河中的气体与尘埃云分布极为不均匀，揭示了星系结构与动力学之间更深层次的联系。
\subsection{早年生活与教育}
钱德拉塞卡于**1910年10月19日**出生在**英属印度的拉合尔（Lahore，今属巴基斯坦）**，出身于一个**泰米尔人家庭**。[11] 他的父母分别是**希塔·巴拉克里希南（Sita Balakrishnan，1891–1931）**与**钱德拉塞卡拉·苏布拉马尼亚·艾耶尔（Chandrasekhara Subrahmanya Ayyar，1885–1960）**。[12] 钱德拉塞卡出生时，其父在**英属印度西北铁路局（Northwestern Railways）**担任**副审计长（Deputy Auditor General）**，因此一家居住在拉合尔。钱德拉塞卡（常被称为 **Chandra**）有两位姐姐——**拉贾拉克希米（Rajalakshmi）**与**巴拉帕尔瓦蒂（Balaparvathi）**；三位弟弟——**维什瓦纳坦（Vishwanathan）**、**巴拉克里希南（Balakrishnan）**与**拉马纳坦（Ramanathan）**；以及四位妹妹——**萨拉达（Sarada）**、**维迪娅（Vidya）**、**萨维特丽（Savitri）**和**森达丽（Sundari）**。他的**叔父**正是印度著名物理学家、诺贝尔奖得主**钱德拉塞卡拉·文卡塔·拉曼（Chandrasekhara Venkata Raman）**。他的母亲是一位热爱学术与文学的女性，曾将**易卜生（Henrik Ibsen）**的戏剧《**玩偶之家（A Doll’s House）**》翻译成泰米尔语，被认为是**启发钱德拉塞卡早年智识兴趣**的重要人物。[13]1916年，钱德拉塞卡一家从拉合尔迁往**阿拉哈巴德（Allahabad）**，并于**1918年定居在马德拉斯（Madras，今金奈）**。

钱德拉塞卡在**12岁之前一直由家中教师教授课程**。[13] 在中学阶段，他的父亲教授他**数学与物理**，而母亲则教他**泰米尔语**。1922年至1925年间，他就读于**马德拉斯特里普利坎的印度中学（Hindu High School, Triplicane, Madras）**。随后，他进入**马德拉斯总统学院（Presidency College, Madras）**就读，该校隶属于**马德拉斯大学（University of Madras）**。1925年至1930年间，他在此完成本科课程。1929年，他受到**阿诺德·索末菲（Arnold Sommerfeld）**一次讲座的启发，撰写了自己的第一篇论文——《**康普顿散射与新统计（The Compton Scattering and the New Statistics）**》。[14]1930年6月，他以优异成绩获得**物理学荣誉学士学位（BSc Hons. in Physics）**。同年7月，钱德拉塞卡获得**印度政府奖学金（Government of India Scholarship）**，前往**剑桥大学（University of Cambridge）**攻读研究生，并被**三一学院（Trinity College）**录取——这一机会由他曾通信交流的物理学家**R. H. 福勒（R. H. Fowler）**促成。在赴英旅途中，钱德拉塞卡利用航行时间推导出**白矮星中简并电子气体的统计力学理论**，并在福勒此前研究的基础上，加入了**相对论修正项**（详见后文“学术遗产”部分）。
\subsection{剑桥大学时期}
在进入**剑桥大学**的第一年，作为**R. H. 福勒（R. H. Fowler）**的研究生，钱德拉塞卡主要从事**平均不透明度（mean opacities）**的计算，并将所得结果应用于**改进简并恒星极限质量模型**的构建。在**英国皇家天文学会（Royal Astronomical Society）**的会议上，他结识了著名天体物理学家**E. A. 米尔恩（E. A. Milne）**。

1931年夏天（研究生第二年），他受**马克斯·玻恩（Max Born）**邀请前往**哥廷根大学（University of Göttingen）**的理论物理研究所工作，研究主题包括**不透明度、原子吸收系数**及**恒星光球层模型**。在**保罗·狄拉克（Paul Dirac）**的建议下，他在最后一年（1932年9月至1933年5月）前往**哥本哈根理论物理研究所（Institute for Theoretical Physics in Copenhagen）**深造，在那里他结识了**尼尔斯·玻尔（Niels Bohr）**。在哥本哈根期间，他与**维克托·魏斯科普夫（Victor Weisskopf）**、**莱昂·罗森菲尔德（Léon Rosenfeld）**、**乔治·普拉泽克（George Placzek）**及**马克斯·德尔布吕克（Max Delbrück）**等人成为好友——他们当时都住在同一家旅馆。

因在**简并恒星研究**方面的成果，钱德拉塞卡获得了**铜质奖章**，并于1933年夏季获得剑桥大学**哲学博士（PhD）**学位，其论文题为《**旋转自引力多项式星体（Rotating Self-Gravitating Polytropes）**》。1933年10月9日，他被选为**剑桥大学三一学院奖学金研究员（Prize Fellow of Trinity College）**，任期为1933年至1937年，成为继**斯里尼瓦萨·拉马努金（Srinivasa Ramanujan）**之后，第二位获得三一学院奖学金的印度人（间隔16年）。有趣的是，他原本**确信自己无法获得该奖学金**，因此已提前计划赴**牛津大学（Oxford）**在米尔恩门下继续研究，并甚至**在当地租好了公寓**。[14]

1934年，钱德拉塞卡访问了**苏联列宁格勒（今圣彼得堡）**，结识了多位著名天体物理学家，包括**维克托·安巴楚缅（Viktor Ambartsumian）**与**尼古拉·亚历山德罗维奇·科济列夫（Nikolai Aleksandrovich Kozyrev）**，并与**列夫·朗道（Lev Landau）**相识。

在此期间，钱德拉塞卡也与英国著名天体物理学家**阿瑟·爱丁顿爵士（Sir Arthur Eddington）**相识。爱丁顿起初对他的研究表现出兴趣，但在**1935年1月**的一次公开演讲中，却**严厉批评了钱德拉塞卡的理论**（详见下节“与爱丁顿的争论”与“钱德拉塞卡—爱丁顿之争”）。
\subsection{职业生涯与研究}
\subsubsection{早期职业生涯}
\begin{figure}[ht]
\centering
\includegraphics[width=6cm]{./figures/8f98ef5906b6192b.png}
\caption{} \label{fig_Sbl_2}
\end{figure}
1935年，钱德拉塞卡应**哈佛天文台（Harvard Observatory）台长哈洛·沙普利（Harlow Shapley）**之邀，前往美国担任为期三个月的**理论天体物理学访问讲师**。他于当年12月启程赴美。在哈佛期间，钱德拉塞卡以其卓越的学识与才华给沙普利留下了深刻印象，但他婉拒了对方提供的**哈佛研究奖学金**职位。此时，钱德拉塞卡结识了著名的荷兰天体物理观测学家**赫拉德·柯伊伯（Gerard Kuiper）**，后者是当时**白矮星领域的权威**。柯伊伯不久前刚被**奥托·斯特鲁维（Otto Struve）**招募至**耶基斯天文台（Yerkes Observatory）**，该台隶属于**芝加哥大学（University of Chicago）**，位于威斯康星州威廉斯湾，由该校校长**罗伯特·梅纳德·哈钦斯（Robert Maynard Hutchins）**主管。斯特鲁维早已听闻钱德拉塞卡的学术声誉，正考虑聘请他担任三位天体物理学教员之一，另两位候选人分别是柯伊伯与丹麦理论物理学家**本特·斯特龙格伦（Bengt Strömgren）**。[14]在柯伊伯的推荐下，斯特鲁维于**1936年3月**邀请钱德拉塞卡前往耶基斯天文台，并向他正式提供职位。钱德拉塞卡对此极感兴趣，但最初仍婉拒邀请，返回英国。不久之后，校长哈钦斯在他返航途中特地**发去电报**敦促他接受职位。最终，钱德拉塞卡同意，并于**1936年12月**重返美国，出任**芝加哥大学理论天体物理学助理教授**。[14]此外，哈钦斯还曾在一起种族歧视事件中出面干预。当时，院长**亨利·盖尔（Henry Gale）**以种族理由否决了钱德拉塞卡参与**斯特鲁维组织的课程教学**的资格。对此，哈钦斯坚定地回应道：“无论如何，让钱德拉塞卡先生去教这门课。”[15]

钱德拉塞卡的整个职业生涯都在**芝加哥大学（University of Chicago）**度过。1941年，他被晋升为**副教授**，两年后的1943年，年仅33岁的他便升任为**正教授**。[14]1946年，**普林斯顿大学（Princeton University）**向钱德拉塞卡提供了一个职位，以接替**亨利·诺里斯·拉塞尔（Henry Norris Russell）**的空缺，并开出**芝加哥大学薪资的两倍**。为了挽留他，校长**罗伯特·哈钦斯（Robert Hutchins）**立即将他的薪资提高至与普林斯顿相同的水平，并成功说服他留在芝加哥。1952年，经**恩里科·费米（Enrico Fermi）**邀请，钱德拉塞卡被任命为**莫顿·D·赫尔理论天体物理学杰出服务教授（Morton D. Hull Distinguished Service Professor of Theoretical Astrophysics）**，并加入**恩里科·费米研究所（Enrico Fermi Institute）**。
1953年，他与妻子**拉莉莎·钱德拉塞卡（Lalitha Chandrasekhar）**正式**加入美国国籍**。[16]

1966年，美国**国家航空航天局（NASA）**在芝加哥大学建立了**天体物理与空间研究实验室（Laboratory for Astrophysics and Space Research，LASR）**。钱德拉塞卡在该大楼二层四个角落办公室之一办公（其余三位入驻者为**约翰·A·辛普森（John A. Simpson）**、**彼得·迈耶（Peter Meyer）**与**尤金·N·帕克（Eugene N. Parker）**）。20世纪60年代末，高层公寓建成后，钱德拉塞卡居住于**莱克肖尔大道4800号（4800 Lake Shore Drive）**，晚年则搬至多彻斯特大楼5550号。
\subsubsection{与爱丁顿的争论}
从剑桥大学毕业后，钱德拉塞卡与**阿瑟·爱丁顿（Arthur Eddington）**保持着密切联系。1935年，他在**英国皇家天文学会（Royal Astronomical Society）**会议上发表了关于恒星结构方程的完整解。然而，爱丁顿在他报告之后立即安排了一场演讲，并在会上**公开批评了钱德拉塞卡的理论**。这一事件使钱德拉塞卡深感沮丧，也由此引发了一场持续多年的科学争论。
爱丁顿拒绝接受恒星存在质量上限的观点，并提出了自己的**替代理论模型**。[17]

钱德拉塞卡随后寻求多位著名物理学家的支持，包括**莱昂·罗森菲尔德（Léon Rosenfeld）**、**尼尔斯·玻尔（Niels Bohr）**与**克里斯蒂安·穆勒（Christian Møller）**等人，他们普遍认为**爱丁顿的反驳缺乏科学依据**。在整个1930年代，两人之间的紧张关系持续存在。爱丁顿在多个学术会议上**公开批评钱德拉塞卡的研究**，双方也多次在论文中互相对比与辩驳。最终，钱德拉塞卡于**1939年**完成了他关于**白矮星理论（Theory of White Dwarfs）**的系统研究，并得到了学界广泛的赞誉。爱丁顿于**1944年去世**，而尽管两人曾有严重分歧，钱德拉塞卡仍多次公开表示，他**深深敬佩爱丁顿**，并一直**视他为朋友**。[17]
\subsubsection{第二次世界大战时期}
在**第二次世界大战**期间，钱德拉塞卡受聘于**马里兰州阿伯丁试验场（Aberdeen Proving Ground）**的**弹道研究实验室（Ballistic Research Laboratory）**工作。期间，他研究了多项与**弹道学和冲击波物理**相关的问题，并撰写了多份技术报告，包括：《**平面激波的衰减（On the Decay of Plane Shock Waves）**》（1943年）《**105毫米炮弹的最佳爆炸高度（Optimum Height for the Bursting of a 105mm Shell）**》《**三重激波存在的条件（On the Conditions for the Existence of Three Shock Waves）**》[18]《**由弹丸反射波干涉所产生拍频波确定弹丸速度的方法》[19]《**爆炸波的正常反射（The Normal Reflection of a Blast Wave）**》[20][9]由于他在**流体力学与气体动力学（hydrodynamics）**方面的专长，**罗伯特·奥本海默（Robert Oppenheimer）**曾邀请他加入**洛斯阿拉莫斯（Los Alamos）**的**曼哈顿计划（Manhattan Project）**。
然而，由于**安全许可（security clearance）审批延误**，钱德拉塞卡最终未能参与该计划。据传，他当时曾**短暂访问过加洛川电磁分离项目（Calutron Project）**。
\subsubsection{系统化的研究哲学}
钱德拉塞卡曾写道，他从事科学研究的动力源于希望“**以自己最大的能力参与科学各个领域的进步**”。在他看来，贯穿其全部学术工作的**首要动机是系统化（systematization）**。
他在论文中如此阐述：“科学家所努力追求的，本质上是选择一个特定的领域、一个特定的方面或一个特定的细节，观察它能否在一个具有形式与连贯性的整体体系中找到恰当的位置；如果不能，他就必须继续探索更多信息，以帮助他实现这一目标。”[21]

钱德拉塞卡形成了一种**独特的研究风格**——他善于在多个物理与天体物理学领域中**逐一精通、逐一系统化**。因此，他的科研生涯可被划分为若干**相对独立的阶段**。
他在每个阶段中都会**深入研究某一主题，连续发表多篇论文**，随后撰写一本**系统总结该领域核心思想的专著**，然后转向新的研究方向。这种系统化的研究节奏贯穿了他的一生：1929–1939年**：研究**恒星结构**，尤其是**白矮星理论**；1939–1943年**：转向**恒星动力学**与**布朗运动理论**；1943–1950年**：研究**辐射输运理论**与**氢负离子的量子理论**；
1950–1961年**：集中于**湍流理论**、**流体力学**与**磁流体稳定性**；1960年代**：研究**椭球平衡体的平衡与稳定**以及**广义相对论**；1971–1983年**：深入发展**黑洞的数学理论**；1980年代后期**：专注于**引力波碰撞理论。[9]钱德拉塞卡这种“**逐个征服、逐步系统化**”的科研方式，使他在多个物理学分支中都留下了深远的理论遗产。
\subsubsection{与学生的合作}
钱德拉塞卡非常重视与学生的合作，并为此深感自豪。他曾指出，在他长达**半个世纪（约1930–1980）**的科研生涯中，**合作者的平均年龄始终保持在30岁左右**——这意味着他总是与年轻一代共同探索前沿科学。他严格要求学生在获得博士学位前，必须称呼他为“**钱德拉塞卡教授（Prof. Chandrasekhar）**”；而在博士毕业后，学生与其他同事一样，可直接称呼他为“**钱德拉（Chandra）**”。在**20世纪40年代**任职于**耶基斯天文台（Yerkes Observatory）**期间，钱德拉塞卡每周末都会**往返150英里（约240公里）**驱车前往芝加哥大学授课。
当时听他课的学生中有两位后来成为诺贝尔奖得主——**李政道（Tsung-Dao Lee）**与**杨振宁（Chen-Ning Yang）**——而他们的获奖时间甚至早于钱德拉塞卡本人。著名天体物理学家**卡尔·萨根（Carl Sagan）**曾回忆钱德拉塞卡的课堂氛围：“那些毫无准备、提出无关紧要问题的学生，往往会被他‘当场处决’；但凡是有价值的问题，钱德拉塞卡都会给予**极为认真、细致的回应**。”[22]
\subsubsection{其他学术与出版活动}
从**1952年至1971年**，钱德拉塞卡担任《**天体物理学期刊（The Astrophysical Journal）**》主编。[23]1957年，当**尤金·帕克（Eugene Parker）**提交论文，首次提出“**太阳风（solar wind）**”的理论时，**两位审稿人均予以否定**。
然而，钱德拉塞卡作为主编，经过审阅后认为帕克的工作**在数学上没有错误**，于是**顶住压力，于1958年毅然发表了这篇论文**。[24]这项决定后来被证明具有划时代意义。

1990年至1995年间，钱德拉塞卡致力于一个特别的研究项目——他尝试使用**现代微积分语言与方法**，重新阐释**艾萨克·牛顿爵士（Sir Isaac Newton）**的《**自然哲学的数学原理（Philosophiae Naturalis Principia Mathematica）**》中的几何推理过程。
这一工作最终形成著作《**给普通读者的〈原理〉（Newton’s Principia for the Common Reader）**》，于1995年出版。此外，钱德拉塞卡还在**引力波碰撞理论（collision of gravitational waves）**[25]与**代数特殊扰动（algebraically special perturbations）**[26]等课题上进行了深入研究。
\subsection{个人生活}
\begin{figure}[ht]
\centering
\includegraphics[width=8cm]{./figures/eb6692de5b844125.png}
\caption{钱德拉塞卡与他的兄弟姐妹在印度马德拉斯（现金奈）合影。从左至右坐着的分别是：**巴拉（Bala）**、**萨维特丽（Savitri）**、**钱德拉（Chandra）**、萨拉达（Sarada）**、**拉贾姆（Rajam）**；站立者从左至右依次为：**维迪娅（Vidya）**、**巴拉克里希南（Balakrishnan）**、**维什瓦姆（Vishwam）**、**拉马纳坦（Ramanathan）**与森达丽。} \label{fig_Sbl_3}
\end{figure}
钱德拉塞卡是**C. V. 拉曼（C. V. Raman）**的侄子，后者曾于1930年获得**诺贝尔物理学奖**。

1936年9月，钱德拉塞卡与**拉莉莎·多拉伊斯瓦米（Lalitha Doraiswamy）**结婚。两人最初相识于**马德拉斯总统学院（Presidency College）**，当时同为在校学生。
1953年，他正式成为**美国归化公民**。在人格与风格上，许多同事与学生认为钱德拉塞卡**温和、积极、慷慨、谦逊、细致且乐于讨论**；但也有人形容他在**私人事务上较为内敛、严肃、急躁甚至固执**，[22]对那些**嘲讽或否定他研究成果的人，则极少宽恕**。[27]钱德拉塞卡是一位**素食主义者**。[28]

1995年，钱德拉塞卡在**芝加哥大学医院（University of Chicago Hospital）**因**心脏病发作**去世，享年84岁。他早在1975年就曾经历过一次心脏病发作。[22]
他的妻子**拉莉莎·钱德拉塞卡**比他多活了18年，于**2013年9月2日**逝世，享年**102岁**。[29]她是一位**文学与西方古典音乐的热心研究者**。[27]

有一次在讨论《薄伽梵歌》时，钱德拉塞卡说道：“在发言之前，我想先作一个个人声明，以免接下来的言论被误解——**我认为自己是一个无神论者。**”[30]他在其他多次演讲中也表达过相同立场。**卡梅什瓦尔·C·瓦利（Kameshwar C. Wali）**曾引用他说：“我在任何意义上都不是一个宗教人士；事实上，**我认为自己是个无神论者。**”[31]在**1987年10月6日**于**芝加哥大学**接受**凯文·克里斯丘纳斯（Kevin Krisciunas）**采访时，钱德拉塞卡提到：“当然，（奥托·斯特鲁维）知道我是无神论者，但他从未在我面前提起过这个话题。”[32]出于对丈夫**无神论信仰的尊重**，钱德拉塞卡的妻子在家中**从不陈列她随身带来的神像或宗教饰物**。[33]
\subsection{奖项、荣誉与遗产}
\subsubsection{诺贝尔奖}
1983年，钱德拉塞卡因其在恒星结构与演化中关键物理过程的理论研究**而获得**诺贝尔物理学奖的一半**。虽然他接受了这一荣誉，但他对诺贝尔委员会的引文感到失望——因为奖项仅仅提及了他**早期的研究成果**，而未能体现他**一生的系统性贡献**。他认为这在某种程度上**贬低了自己毕生的科学成就**。该奖项的另一半授予了**威廉·A·福勒。
\subsubsection{其他奖项与荣誉}
\begin{itemize}
\item 英国皇家学会院士（Fellow of the Royal Society, FRS）**（1944年）[1]
\item 美国哲学学会（American Philosophical Society）会员**（1945年）[34]
\item 美国艺术与科学院（American Academy of Arts and Sciences）院士**（1946年）[35]
\item 亨利·诺里斯·罗素讲座奖（Henry Norris Russell Lectureship）**（1949年）[36]
\item 布鲁斯奖章（Bruce Medal）**（1952年）[37]
\item 英国皇家天文学会金质奖章（Gold Medal of the Royal Astronomical Society）**（1953年）[38]
\item 美国国家科学院（U.S. National Academy of Sciences）院士**（1955年）[39]
\item 鲁姆福德奖（Rumford Prize）**，美国艺术与科学院颁发（1957年）[40]
\item 美国国家科学奖章（National Medal of Science, USA）**（1966年）[41]
\item 印度国家勋章“莲花装勋章二级”（Padma Vibhushan）**（1968年）
\item 美国国家科学院德雷珀奖章（Henry Draper Medal of the National Academy of Sciences）**（1971年）[42]
\item 玛丽安·斯摩鲁霍夫斯基奖章（Marian Smoluchowski Medal）**（1973年）
\item 英国皇家学会科普利奖章（Copley Medal of the Royal Society）**（1984年）
\item 戈登·J·莱因奖（Gordon J. Laing Award）**（1989年）
\item 美国成就学院金盘奖（Golden Plate Award of the American Academy of Achievement）**（1990年）[43]
\item 詹斯基讲座（Jansky Lectureship）**，于**美国国家射电天文台（National Radio Astronomy Observatory）**发表
\end{itemize}
\subsubsection{学术遗产}
钱德拉塞卡最著名的科学成就之一是**天体物理学中的“钱德拉塞卡极限（Chandrasekhar limit）”**。该极限给出了**白矮星的最大质量**——约为**1.44倍太阳质量**，也即是**恒星在超新星爆发后坍缩成中子星或黑洞所需超过的最小质量**。这一极限最早由钱德拉塞卡于**1930年**在他**从印度前往英国剑桥攻读研究生的首次航行途中**计算得出。1979年，美国**国家航空航天局（NASA）**将其“四大太空天文台（Great Observatories）”之一命名为**钱德拉X射线天文台（Chandra X-ray Observatory）**，以纪念他的贡献。该命名来自一次全球征名活动，共收到来自**50个州与61个国家的6000份提名**。该天文台于**1999年7月23日**由**哥伦比亚号航天飞机（Space Shuttle Columbia）**发射升空并部署入轨。此外，多个重要科学名词与实体以他命名：钱德拉塞卡数（Chandrasekhar number）**：磁流体力学中的一个重要无量纲数；小行星1958 Chandra**；喜马拉雅钱德拉望远镜（Himalayan Chandra Telescope）**。在《**英国皇家学会院士传记回忆录（Biographical Memoirs of Fellows of the Royal Society of London）**》中，天体物理学家 **R. J. 泰勒（R. J. Tayler）** 曾写道：“钱德拉塞卡是一位典型的应用数学大师——他的研究主要服务于天文学，而像他这样的人，或许再也不会出现了。”[1]

在职业生涯中，钱德拉塞卡共指导了**45名博士生**。[44]他去世后，妻子**拉莉莎·钱德拉塞卡（Lalitha Chandrasekhar）**将他的**诺贝尔奖金全数捐赠给芝加哥大学**，用于设立**“苏布拉马尼扬·钱德拉塞卡纪念奖学金（Subrahmanyan Chandrasekhar Memorial Fellowship）”**。该奖学金于**2000年首次颁发**，每年授予在**物理系或天文学与天体物理学系博士项目**中表现最为杰出的申请者。[45]此外，自**2014年**起，**亚太物理学会联合会（AAPS）**设立了**“钱德拉塞卡等离子体物理奖（S. Chandrasekhar Prize of Plasma Physics）”**，以表彰在**等离子体物理领域**作出杰出贡献的科学家。[46]

**钱德拉天体物理学院（The Chandra Astrophysics Institute，CAI）** 是一个面向**高中生**的教育项目，由**麻省理工学院（MIT）**的科学家担任导师[47]，并由**钱德拉X射线天文台（Chandra X-ray Observatory）**资助[48]。该项目旨在培养青年学生对天体物理学的兴趣与研究能力。著名天体物理学家**卡尔·萨根（Carl Sagan）**在其著作《**魔鬼出没的世界（The Demon-Haunted World）**》中称赞道：“我是在苏布拉马尼扬·钱德拉塞卡那里，第一次领悟到什么是真正的数学优雅。”2017年10月19日，为纪念钱德拉塞卡诞辰107周年及其“**钱德拉塞卡极限（Chandrasekhar limit）**”的发现，**谷歌（Google）**在**28个国家**同步发布了纪念他的**Google Doodle**。[49][50]

2010年，值钱德拉塞卡百年诞辰之际，**芝加哥大学（University of Chicago）**举办了题为**“钱德拉塞卡百年纪念研讨会（Chandrasekhar Centennial Symposium 2010）”**的学术会议。与会的世界顶尖天体物理学家包括：**罗杰·彭罗斯（Roger Penrose）**、**基普·索恩（Kip Thorne）**、**弗里曼·戴森（Freeman Dyson）**、**贾扬特·纳利卡尔（Jayant V. Narlikar）**、**拉希德·苏尼亚耶夫（Rashid Sunyaev）**、**G. 斯里尼瓦桑（G. Srinivasan）**与**克利福德·威尔（Clifford Will）**等人。研讨会的研究报告随后于2011年结集出版，书名为：《**从流体流动到黑洞：献给钱德拉塞卡百年诞辰的致敬（Fluid Flows to Black Holes: A Tribute to S. Chandrasekhar on His Birth Centenary）**》[51][52][53]。
\subsection{出版著作}
\subsubsection{书籍}
\begin{itemize}
\item Chandrasekhar, S.**（2010）[1939/1958]，《**恒星结构研究导论（An Introduction to the Study of Stellar Structure）**》，纽约：多佛出版社（Dover）。ISBN 978-0-486-60413-8。
\item Chandrasekhar, S.**（2005）[1942]，《**恒星动力学原理（Principles of Stellar Dynamics）**》，纽约：多佛出版社。ISBN 978-0-486-44273-0。
\item Chandrasekhar, S.**（2011）[1950/1960]，《**辐射输运（Radiative Transfer）**》，纽约：多佛出版社。ISBN 978-0-486-60590-6。
\item Chandrasekhar, S.**（1975）[1960]，《**等离子体物理学（Plasma Physics）**》，芝加哥：芝加哥大学出版社（The University of Chicago Press）。ISBN 978-0-226-10084-5。
\item Chandrasekhar, S.**（1981）[1961]，《**流体力学与磁流体稳定性（Hydrodynamic and Hydromagnetic Stability）**》，纽约：多佛出版社。ISBN 978-0-486-64071-6。
\item Chandrasekhar, S.**（1987）[1969]，《**椭球平衡体（Ellipsoidal Figures of Equilibrium）**》，纽约：多佛出版社。ISBN 978-0-486-65258-0。
\item Chandrasekhar, S.**（1998）[1983]，《**黑洞的数学理论（The Mathematical Theory of Black Holes）**》，纽约：牛津大学出版社（Oxford University Press）。ISBN 978-0-19-850370-5。
\item Chandrasekhar, S.**（1983）[1983]，《**爱丁顿：他那个时代最杰出的天体物理学家（Eddington: The Most Distinguished Astrophysicist of His Time）**》，剑桥大学出版社（Cambridge University Press）。ISBN 9780521257466。
\item Chandrasekhar, S.**（1990）[1987]，《**真理与美：科学中的审美与动机（Truth and Beauty: Aesthetics and Motivations in Science）**》，芝加哥：芝加哥大学出版社。ISBN 978-0-226-10087-6。
\item Chandrasekhar, S.**（1995），《**给普通读者的〈原理〉（Newton's Principia for the Common Reader）**》，牛津：克拉伦登出版社（Clarendon Press）。ISBN 978-0-19-851744-3。
\item Spiegel, E. A.**（2011）[1954]，《**湍流理论：钱德拉塞卡1954年讲义（The Theory of Turbulence: Subrahmanyan Chandrasekhar’s 1954 Lectures）**》，荷兰：施普林格出版社（Springer）。ISBN 978-94-007-0117-5。
\end{itemize}
\subsubsection{注释}
\begin{itemize}
\item Chandrasekhar, S.**（1939）。〈**恒星系统的动力学 I–VIII（The Dynamics of Stellar Systems. I–VIII）**〉，《天体物理学期刊（The Astrophysical Journal）》，第90卷第1期，页1–154。Bibcode: 1939ApJ....90....1C。DOI: 10.1086/144094。ISSN 0004-637X。
\item Chandrasekhar, S.**（1943）。〈**物理学与天文学中的随机问题（Stochastic Problems in Physics and Astronomy）**〉，《现代物理评论（Reviews of Modern Physics）》，第15卷第1期，页1–89。Bibcode: 1943RvMP...15....1C。DOI: 10.1103/RevModPhys.15.1。ISSN 0034-6861。
\item Chandrasekhar, S.**（1993）。《**经典广义相对论（Classical General Relativity）**》，英国皇家学会（Royal Society）。
\item Chandrasekhar, S.**（1979）。〈**广义相对论的作用：回顾与展望（The Role of General Relativity: Retrospect and Prospect）**〉，国际天文学联合会会议论文集（Proc. IAU Meeting）。[54]
\item Chandrasekhar, S.**（1943）。《**恒星动力学新方法（New Methods in Stellar Dynamics）**》，纽约科学院（New York Academy of Sciences）。
\item Chandrasekhar, S.**（1954）。〈**瑞利散射下阳光天空的照度与偏振（The Illumination and Polarization of the Sunlit Sky on Rayleigh Scattering）**〉，《美国哲学学会汇刊（Transactions of the American Philosophical Society）》，第44卷第6期，页643–728。DOI: 10.2307/1005777。JSTOR 1005777。
\item Chandrasekhar, S.**（1983）。〈**恒星：它们的演化与稳定性（On Stars, Their Evolution and Their Stability）**，诺贝尔奖演讲（Nobel Lecture）〉，《现代物理评论（Reviews of Modern Physics）》，第56卷第2期，页137–147。斯德哥尔摩：诺贝尔基金会。DOI: 10.1103/RevModPhys.56.137。
\item Chandrasekhar, S.**（1981）。《**人类知识的新视野（New Horizons of Human Knowledge）**：在联合国教科文组织的系列公开讲座》，巴黎：联合国教科文组织出版社（Unesco Press）。
\item Chandrasekhar, S.**（1975）。〈**莎士比亚、牛顿与贝多芬：或，创造力的模式（Shakespeare, Newton, and Beethoven: Or, Patterns of Creativity）**〉，《当代科学（Current Science）》，第70卷第9期，页810–822。芝加哥大学出版。JSTOR 24099932。
\item Chandrasekhar, S.**（1973年7月）。〈**P.A.M. 狄拉克七十寿辰纪念（P.A.M. Dirac on His Seventieth Birthday）**〉，《当代物理学（Contemporary Physics）》，第14卷第4期，页389–394。Bibcode: 1973ConPh..14..389C。DOI: 10.1080/00107517308210761。ISSN 0010-7514。
\item Chandrasekhar, S.**（1947）。收录于 **Robert B. Heywood（编）**，《**心灵的工作：科学家篇（The Works of the Mind: The Scientist）**》，芝加哥大学出版社（University of Chicago Press），页159–179。OCLC 752682744。
\item Chandrasekhar, S.**（1995）。《**关于拉马努金雕像的回忆与发现（Reminiscences and Discoveries on Ramanujan’s Bust）**》，英国皇家学会出版。ASIN B001B12NJ8。
\item Chandrasekhar, S.**（1990）。《**如何探索广义相对论的物理内容（How One May Explore the Physical Content of the General Theory of Relativity）**》，美国数学学会（American Mathematical Society）。ASIN B001B10QTM。
\end{itemize}
\subsubsection{期刊论文}
钱德拉塞卡一生共发表了约 **380 篇论文**[55][1]。
他在 **1928 年**，仍为本科生时，发表了第一篇论文，主题是**康普顿效应（Compton Effect）**[56]；而他最后一篇论文发表于 **1995 年**，距他去世仅两个月，被接受发表的主题是**恒星的非径向振荡（Non-radial Oscillation of Stars）**[57]。芝加哥大学出版社（University of Chicago Press）** 将他的重要论文编为七卷本的《钱德拉塞卡精选论文集（Selected Papers of S. Chandrasekhar）》出版。
\begin{itemize}
\item Chandrasekhar, S.**（1989）。《**精选论文，第1卷：恒星结构与恒星大气（Selected Papers, Vol. 1: Stellar Structure and Stellar Atmospheres）**》，芝加哥：芝加哥大学出版社。ISBN 9780226100890。
\item Chandrasekhar, S.**（1989）。《**精选论文，第2卷：辐射输运与氢负离子（Selected Papers, Vol. 2: Radiative Transfer and Negative Ion of Hydrogen）**》，芝加哥：芝加哥大学出版社。ISBN 9780226100920。
\item Chandrasekhar, S.**（1989）。《**精选论文，第3卷：物理与天文学中的随机、统计与磁流体问题（Selected Papers, Vol. 3: Stochastic, Statistical and Hydromagnetic Problems in Physics and Astronomy）**》，芝加哥：芝加哥大学出版社。ISBN 9780226100944。
\item Chandrasekhar, S.**（1989）。《**精选论文，第4卷：等离子体物理学、流体与磁流体稳定性及张量-维里定理的应用（Selected Papers, Vol. 4: Plasma Physics, Hydrodynamic and Hydromagnetic Stability, and Applications of the Tensor-Virial Theorem）**》，芝加哥：芝加哥大学出版社。ISBN 9780226100975。
\item Chandrasekhar, S.**（1990）。《**精选论文，第5卷：相对论天体物理学（Selected Papers, Vol. 5: Relativistic Astrophysics）**》，芝加哥：芝加哥大学出版社。ISBN 9780226100982。
\item Chandrasekhar, S.**（1991）。《**精选论文，第6卷：黑洞与平面引力波碰撞的数学理论（Selected Papers, Vol. 6: The Mathematical Theory of Black Holes and of Colliding Plane Waves）**》，芝加哥：芝加哥大学出版社。ISBN 9780226101019。
\item Chandrasekhar, S.**（1997）。《**精选论文，第7卷：广义相对论中的恒星非径向振荡及其他著作（Selected Papers, Vol. 7: The Non-radial Oscillations of Stars in General Relativity and Other Writings）**》，芝加哥：芝加哥大学出版社。ISBN 9780226101040。
\end{itemize}
\subsubsection{关于钱德拉塞卡的书籍与论文}
\begin{itemize}
\item Miller, Arthur I.**（2005）。《**群星帝国：追寻黑洞之旅中的友情、执念与背叛（Empire of the Stars: Friendship, Obsession, and Betrayal in the Quest for Black Holes）**》，波士顿：霍顿·米夫林出版社（Houghton Mifflin）。ISBN 978-0-618-34151-1。

\item Srinivasan, G.（编）**（1997）。《**从白矮星到黑洞：钱德拉塞卡的遗产（From White Dwarfs to Black Holes: The Legacy of S. Chandrasekhar）**》，芝加哥：芝加哥大学出版社（The University of Chicago Press）。ISBN 978-0-226-76996-7。

\item Penrose, Roger**（1996）。〈**钱德拉塞卡、黑洞与奇点（Chandrasekhar, Black Holes and Singularities）**〉，《天体物理与天文学杂志（Journal of Astrophysics and Astronomy）》，第17卷第3–4期，页213–231。Bibcode: 1996JApA...17..213P。DOI: 10.1007/BF02702305。ISSN 0250-6335。

\item Parker, E.**（1996）。〈**钱德拉塞卡与磁流体力学（S. Chandrasekhar and Magnetohydrodynamics）**〉，《天体物理与天文学杂志》，第17卷第3–4期，页147–166。DOI: 10.1007/BF02702301。ISSN 0250-6335。

\item Wali, Kameshwar C.**（1991）。《**钱德拉：S. 钱德拉塞卡传（Chandra: A Biography of S. Chandrasekhar）**》，芝加哥：芝加哥大学出版社。ISBN 978-0-226-87054-0。

\item Wali, Kameshwar C.（编）**（1997）。《**钱德拉塞卡：传奇背后的男人——怀念钱德拉（Chandrasekhar: The Man Behind the Legend – Chandra Remembered）**》，伦敦：帝国理工学院出版社（Imperial College Press）。ISBN 978-1-86094-038-5。

\item Wali, Kameshwar C.（编）**（2001）。《**视角的追寻（A Quest for Perspectives）**》，新加坡：世界科学出版社（World Scientific Publishing Co. Pte. Ltd.）。ISBN 978-1-86094-201-3。

\item Wali, Kameshwar C.（编）**（2020）。《**S. 钱德拉塞卡：精选通信与对话（S. Chandrasekhar: Selected Correspondence and Conversations）**》，世界科学出版社。ISBN 978-9811208324。

\item Wignesan, T.（编）**（2004）。〈**让群星黯然失色的人（The Man Who Dwarfed the Stars）**〉，《亚洲学者的亚洲（The Asianists’ Asia）》。ISSN 1298-0358。

\item Venkataraman, G.**（1992）。《**钱德拉塞卡与他的极限（Chandrasekhar and His Limit）**》，印度海得拉巴：大学出版社（Universities Press）。ISBN 978-81-7371-035-3。

\item Saikia, D. J. 等（编）**（2011）。《**流体流向黑洞：钱德拉塞卡百年诞辰纪念论文集（Fluid Flows to Black Holes: A Tribute to S. Chandrasekhar on His Birth Centenary）**》，新加坡：世界科学出版社。ISBN 978-981-4299-57-2。

\item Ramnath, Radhika（编）**（2012）。《**S. 钱德拉塞卡：科学之人（S. Chandrasekhar: Man of Science）**》，哈珀柯林斯出版社（HarperCollins）。ASIN B00C3EWIME。

\item Alic, Kameshwar C.；Wali, Kameshwar C.（编）**（2011）。《**科学自传：S. 钱德拉塞卡（A Scientific Autobiography: S. Chandrasekhar）**》，世界科学出版社。DOI: 10.1142/7686。ISBN 978-981-4299-57-2。

\item Salwi, Dilip（编）**（2004）。《**S. 钱德拉塞卡：学者型科学家（S. Chandrasekhar: The Scholar Scientist）**》，Rupa出版社。ISBN 978-8129104915。

\item Pandey, Rakesh Kumar（编）**（2017）。《**钱德拉塞卡极限：白矮星的大小（Chandrasekhar Limit: Size of White Dwarfs）**》，Lap Lambert Academic Publishing。ISBN 978-3330317666。

\end{itemize}